% © 2026 Ing. David Jaroš — CC BY-NC-ND 4.0
%
% This work is licensed under a Creative Commons Attribution-NonCommercial-NoDerivatives
% 4.0 International License (CC BY-NC-ND 4.0).
%
% License History: Earlier drafts (up to v0.3) were released under CC BY 4.0.
% From v0.4 onward, all material is released under CC BY-NC-ND 4.0 to protect
% the integrity of the theoretical work during ongoing academic development.

\documentclass[11pt,a4paper]{article}
\usepackage{amsmath, amssymb, amsthm}
\usepackage{hyperref}
\usepackage{geometry}
\geometry{margin=2.5cm}

\newtheorem{theorem}{Theorem}
\newtheorem{proposition}{Proposition}
\newtheorem{definition}{Definition}
\newtheorem{remark}{Remark}
\newtheorem{conjecture}{Conjecture}

\title{$B_{\mathrm{base}}$ from Entropy and Spectral Counting\\
\large Attempt E-1: Is the Exponent $\tfrac{3}{2}$ Entropic?}
\author{Ing.\ David Jaroš}
\date{2026}

\begin{document}
\maketitle

\begin{abstract}
We investigate whether the exponent $\tfrac{3}{2}$ in the UBT formula
$B_{\mathrm{base}} = N_{\mathrm{eff}}^{3/2}$ arises from entropy or spectral
counting of $\Theta$-modes.  Three candidate mechanisms are explored:
(1) the thermal partition function on the torus; (2) the counting exponent
from the density of states of the KK tower; and (3) the information entropy
of the KK mode distribution.  All approaches reach well-defined conclusions
about their status.  This is attempt E-1 in the $B_{\mathrm{base}}$ programme;
the full list is in \texttt{research/alpha\_derivation\_status.md}.
All claims are labelled with their proof status.
\end{abstract}

\tableofcontents

% ======================================================================
\section{Setup: The $B_{\mathrm{base}}$ Problem}
\label{sec:setup}
% ======================================================================

\subsection{The formula}

Following \texttt{research/alpha\_derivation\_status.md} and
\texttt{canonical/interactions/qed.tex} §4--5, the UBT fine-structure constant
formula is:
\begin{equation}
\label{eq:alpha_formula}
  \alpha = \frac{B_0}{B_{\mathrm{base}} + B_\alpha},
  \quad B_0 = 8\pi, \quad B_\alpha \approx 46.3,
  \quad B_{\mathrm{base}} \approx 41.57 \approx N_{\mathrm{eff}}^{3/2},
\end{equation}
where $N_{\mathrm{eff}} \approx 11.53$ is motivated by BRST cohomology
(Assumptions A1--A3, \texttt{THEORY/axioms/core\_assumptions.tex}).

\textbf{Open problem:} Why is the exponent exactly $\tfrac{3}{2}$?

% ======================================================================
\section{Entropy Approach: Thermal Partition Function on the Torus}
\label{sec:thermal}
% ======================================================================

\subsection{Setup}

Consider the $\Theta$-field at finite temperature $T = 1/\beta$ on the
spacetime torus $\mathbb{T}^3 \times S^1_\beta$.  The thermal partition
function is:
\begin{equation}
\label{eq:Z_thermal}
  Z(\beta) = \mathrm{Tr}\, e^{-\beta H}
  = \prod_{n \in \mathbb{Z}} \frac{1}{1 - e^{-\beta m_n}},
\end{equation}
where $m_n = |n|/R_\psi$ are the KK masses.

\subsection{High-temperature expansion}

In the high-temperature limit $\beta m_1 = \beta/R_\psi \ll 1$, the
partition function is dominated by the entropy $S = \beta^2 \partial_\beta \ln Z$:
\begin{equation}
\label{eq:entropy_ht}
  S \;\sim\; \frac{\pi^2}{3}\, \frac{N_{\mathrm{dof}}}{\beta/R_\psi},
\end{equation}
where $N_{\mathrm{dof}}$ is the number of light degrees of freedom.

\begin{remark}[Does this give $N^{3/2}$?]
  The high-temperature entropy \eqref{eq:entropy_ht} scales as
  $S \sim N_{\mathrm{dof}} (T R_\psi)$, which is linear in $N_{\mathrm{dof}}$,
  not $N_{\mathrm{dof}}^{3/2}$.  The thermal approach in this simple form
  does \emph{not} produce the $\tfrac{3}{2}$ exponent.

  Status: \textbf{[Dead end E1-a]}.
\end{remark}

% ======================================================================
\section{Density-of-States Approach}
\label{sec:dos}
% ======================================================================

\subsection{Spectral counting}

For a field on a $d$-dimensional torus of volume $V_d$, the number of modes
with frequency below $\omega$ is (Weyl's law):
\begin{equation}
\label{eq:weyl}
  N(\omega) = \frac{V_d\,\omega^d}{(2\pi)^d\,\Gamma(d/2+1)}.
\end{equation}

For $d = 3$ compact dimensions:
\begin{equation}
  N(\omega) = \frac{V_3\,\omega^3}{6\pi^2}.
\end{equation}

\subsection{Free energy scaling}

The free energy of the mode sum scales as:
\begin{equation}
\label{eq:F_dos}
  F \;\sim\; -\int_0^\Lambda d\omega\; N(\omega)\, \omega
  \;\sim\; -V_3 \Lambda^4 / (8\pi^2).
\end{equation}

\begin{remark}[Does this give $N^{3/2}$?]
  The free energy in \eqref{eq:F_dos} scales as $\Lambda^4$, not $N^{3/2}$.
  To obtain $N^{3/2}$ from a density-of-states argument, one would need
  a two-dimensional compact space ($d=2$) with an additional half-power
  from a Casimir-like contribution.  This possibility is not excluded but
  requires a specific geometric setup that has not been identified.

  Status: \textbf{[Inconclusive E1-b]}.
\end{remark}

% ======================================================================
\section{Information Entropy of the KK Distribution}
\label{sec:info}
% ======================================================================

\subsection{Von Neumann entropy of the mode density matrix}

Define the normalised KK mode density:
\begin{equation}
  p_n = \frac{e^{-\beta m_n}}{Z(\beta)}, \quad n \in \mathbb{Z}.
\end{equation}

The von Neumann entropy is:
\begin{equation}
\label{eq:vn_entropy}
  S_{\mathrm{vN}} = -\sum_{n} p_n \ln p_n.
\end{equation}

\begin{proposition}[Scaling of KK entropy]\label{prop:kk_entropy}
  For $N_{\mathrm{modes}}$ modes with equal weights $p_n = 1/N_{\mathrm{modes}}$:
  \begin{equation}
    S_{\mathrm{vN}} = \ln N_{\mathrm{modes}}.
  \end{equation}
  This grows as $\ln N$, not $N^{3/2}$.
\end{proposition}

\begin{remark}[Does this give $N^{3/2}$?]
  The von Neumann entropy scales logarithmically in $N$, not as $N^{3/2}$.
  One could attempt to identify $B_{\mathrm{base}} = e^{S_{\mathrm{vN}}}$
  (exponential of entropy), which gives $B_{\mathrm{base}} = N$, not $N^{3/2}$.
  To obtain $N^{3/2}$, one would need a \emph{Rényi entropy} of order
  $\alpha_R = -1/2$, which has no clear physical motivation.

  Status: \textbf{[Dead end E1-c]}.
\end{remark}

% ======================================================================
\section{Partial Positive Result: Boltzmann Counting on $S^2$}
\label{sec:s2}
% ======================================================================

\subsection{Mode counting on a two-sphere}

If the compact imaginary-time directions $(\psi, \chi, \xi)$ form a two-sphere
$S^2$ (rather than three independent circles), then the spherical harmonic
spectrum has degeneracy $2\ell+1$ at level $\ell$, and the total mode count
up to level $L$ is:
\begin{equation}
  N_{\leq L} = \sum_{\ell=0}^L (2\ell+1) = (L+1)^2.
\end{equation}

The partition function at temperature $T$ on $S^2$ (radius $R$) is:
\begin{equation}
  Z_{S^2}(T) = \prod_{\ell=0}^{\infty} \prod_{m=-\ell}^{\ell}
  \frac{1}{1 - e^{-m_{\ell}/T}}, \quad
  m_\ell = \sqrt{\ell(\ell+1)}/R.
\end{equation}

At high temperature, this produces:
\begin{equation}
  \ln Z_{S^2} \;\sim\; (TR)^2 \cdot \zeta(3),
\end{equation}
and the free energy scales as $F \sim -T^3 R^2$, giving
$B \sim (N_{\mathrm{eff}} \cdot R^2)^{3/2}$ schematically.

\begin{remark}[Partial success: geometry-dependent]
  The $S^2$ mode counting produces a term with the correct $N^{3/2}$ scaling
  \emph{if} $N_{\mathrm{eff}} \propto (TR)^2$.  However:
  \begin{enumerate}
    \item The compact imaginary directions in UBT are three circles
          $S^1 \times S^1 \times S^1$, not $S^2$;
    \item The specific numerical coefficient does not match $B_{\mathrm{base}} = 41.57$;
    \item The temperature scale is undetermined.
  \end{enumerate}
  Status: \textbf{[L2]} — a partial structural observation, not a derivation.
\end{remark}

% ======================================================================
\section{Conclusions and Next Steps}
\label{sec:conclusions}
% ======================================================================

\begin{center}
\begin{tabular}{lll}
\hline
Approach & Outcome & Next \\
\hline
Thermal entropy (torus) & Linear in $N$ — dead end & E1-a \\
Weyl density of states & $\Lambda^4$ scaling — inconclusive & E1-b \\
Von Neumann entropy & $\ln N$ scaling — dead end & E1-c \\
Boltzmann counting on $S^2$ & $N^{3/2}$ structurally possible & E1-d (open) \\
\hline
\end{tabular}
\end{center}

The most promising direction from this investigation is the $S^2$ geometry
scenario (§\ref{sec:s2}).  The next step is to determine whether the UBT
compact imaginary-time manifold can be globally $S^2$ (rather than
$S^1 \times S^1 \times S^1$) under some dynamical constraint, and whether
this produces the correct numerical coefficient.

This investigation is continued in:
\begin{itemize}
  \item \texttt{research/B\_base\_spectral\_determinant.tex} — spectral
    determinant approach;
  \item \texttt{research/B\_base\_heat\_kernel.tex} — heat kernel analysis.
\end{itemize}

\end{document}
